% SPDX-License-Identifier: Apache-2.0
% covers: LineBuf
\datasheet{LineBuf}{A line of pixels on the clock that shows them,
filled from a channel and read by column.}

\dsfacts{\code{lib/parts/src/dma.rs}}{\code{txhdl\_parts::dma::LineBuf}}%
{\code{//docs:examples}}

\subsection*{Function}

The far end of a scanout. \code{LineFetch} reads a line from memory
on the bus clock and \code{ChanCdc} carries the words to the pixel
clock; this is what they arrive in. A word taken lands at the next
place in the line, and the raster asks for a column and gets what is
there.

\subsection*{Parameters}

\begin{center}\footnotesize
\begin{tabular}{@{}lp{0.68\textwidth}@{}}
\toprule
Parameter & Meaning \\
\midrule
\code{LEN} & How many pixels a line holds. \\
\code{AW} & The width of a column index, in bits. \\
\code{C} & The clock it runs on, which is the pixel clock. \\
\bottomrule
\end{tabular}
\end{center}

The tables below use \code{LineBuf<32, 5, ClkPix>}, the buffer of
\code{ex\_scanout}.

\dstables{LineBuf}

\subsection*{Behaviour}

\begin{itemize}
\item A word is taken whenever one is offered. The buffer is the
consumer that the fetch is held up by, so never refusing is what
keeps the line filling.
\item \code{shown} holds the pixel the column named at the last edge,
and \code{pix} is that register.
\item \code{sol} puts the write back to the start of the line, so a
fetch that delivered too few words does not walk the buffer out of
step for ever.
\end{itemize}

\subsection*{Why the output is registered}

A block RAM has an output register, and a design that reads the
memory straight out to a port asks it to settle combinationally
within a pixel. A raster also wants a pixel that is stable for the
whole cycle it is shown rather than one that moves when the column
does. Both say the same thing, so the port is a register and not the
memory.

It costs one cycle of lead, which a raster has, since it knows the
column it will want next.

\subsection*{Verification}

\begin{itemize}
\item \code{ex\_scanout} joins \code{LineFetch}, \code{ChanCdc} and
this over two clocks of period four and period six, and reads
thirty-two pixels back a column at a time. The netlist is checked
against that run under nvc and Verilator:
\code{//docs:scanout\_sim\_linebuf\_tb\_test} and
\code{//docs:scanout\_vsim\_test}.
\item The memory it reads from holds its own address in each word, so
a pixel that came from the wrong burst, crossed badly, or landed out
of step says which place it came from.
\end{itemize}

\subsection*{Limits}

Writing and reading are independent, and nothing here says the fill
is far enough ahead of the beam. What "far enough" means belongs to
the video timing rather than to a buffer. A line that is not filled
in time shows the previous line's pixel at that column, which is a
visible artefact rather than a hang.

One word a pixel, which is what a framebuffer holds and what the
crossing carries. It does not pack several pixels to a word.

\code{LEN} and \code{AW} are not checked against each other.
